% =======================================================
\section{Introduction}
% =======================================================

Indoor positioning has become a key enabling technology for applications such as robotics, asset tracking, navigation, human--machine interaction, and context-aware wireless services. In this context, visible light positioning (VLP) has emerged as a promising alternative to radio-frequency-based indoor localization because it can exploit existing lighting infrastructure while offering high accuracy, low electromagnetic interference, and favorable integration with optical wireless communication systems \cite{Wang2024,Zhu2025}. Recent surveys further show that received-signal-strength (RSS)-based VLP remains particularly attractive due to its low hardware complexity and compatibility with photodiode receivers, making it relevant for practical low-cost deployments \cite{Wang2024,Zhu2025}.

Most RSS-based VLP architectures, however, rely on multiple fixed LEDs mounted on the ceiling, and estimate the receiver position by means of multilateration or related geometric reconstructions \cite{Wang2024,Almadani2021}. Although effective, this conventional multi-LED paradigm increases infrastructure density and reduces flexibility, especially when moving from 2D positioning to full 3D localization. Experimental studies have also shown that practical 3D VLP performance can be significantly affected by receiver tilt, multipath propagation, and mismatches between the assumed and actual optical behavior of the emitters \cite{Almadani2021}. Other demonstrations of 3D VLP have addressed the problem by increasing receiver complexity, for example through multiple photodiodes or learning-based processing \cite{Zhang2020}, while some single-source approaches require additional beacons or alternative sensing configurations \cite{Huaping2020}. As a result, the experimental validation of 3D VLP with both a minimal receiver and a minimal transmitter infrastructure remains largely open.

At the same time, machine-learning-based VLP has attracted growing interest as a means to compensate for nonlinearities, reflections, and model mismatch \cite{Tran2022}. Nevertheless, such approaches are inherently data-driven and often behave as black-box estimators, whose predictions depend on training data, scenario-specific fingerprints, and implicit nonlinear mappings that are difficult to interpret physically \cite{Tran2022}. By contrast, model-based methods preserve a direct link between the optical channel, the measurement model, and the estimated geometry. This interpretability is especially valuable in optical positioning, where the impact of steering angles, receiver orientation, and RSS uncertainty can be explicitly analyzed and related to the underlying propagation model.

Motivated by this, our recent work in \cite{acuna2025model} introduced a ratio-based 3D positioning framework using a single beam-steered LED (SBL) and a single photodiode receiver. Under ideal Lambertian assumptions, the method cancels the unknown receiver-orientation term through inter-orientation RSS ratios and yields closed-form generalized least-squares (GLS) and weighted least-squares (WLS) estimators. A related 2D study in \cite{acuna2026robust} experimentally demonstrated the intrinsic robustness of the ratio construction against receiver tilt. However, despite these promising results, the experimental feasibility of the full 3D SBL concept has not yet been established.

This paper fills that gap by presenting, to the best of our knowledge, the first experimental validation of 3D VLP based on a single beam-steered LED and a single photodiode receiver. Beyond full 3D localization, the paper also evaluates direction-finding performance, which is directly relevant to beam-steering operation, and examines the effect of integration time on estimation stability. In this way, the paper validates the practical viability of the SBL principle while preserving the core advantage of a model-based, interpretable, non-black-box localization framework.

The remainder of this paper is organized as follows. Section~\ref{sec:system_model} summarizes the SBL system model and the GLS/WLS estimators. Section~\ref{sec:testbed} describes the 3D testbed and the measurement protocol. Section~\ref{sec:results} presents the experimental validation results, the comparison between estimators, the direction-finding analysis, and the position-estimation stability study. Finally, Section~\ref{sec:conclusions} concludes the paper.


% =======================================================
\section{Introduction}
% =======================================================

\IEEEPARstart{I}{ndoor} localization is a key enabler for robotics, navigation, asset tracking, and context-aware wireless services, yet global navigation satellite systems (GNSS) remain unreliable indoors due to severe attenuation, blockage, and multipath. Consequently, a wide range of indoor positioning systems (IPS) based on radio-frequency technologies has been studied, but accuracy in cluttered environments is often limited by non-line-of-sight conditions, time-varying interference, and deployment/calibration constraints \cite{Zafari2019Survey}. These limitations have motivated complementary modalities capable of delivering higher accuracy with practical infrastructure.

In this context, optical wireless positioning (OWP), including visible-light positioning (VLP), leverages LED transmitters and photodiode/camera receivers to transform lighting into a positioning network. Surveyed VLC/VLP systems offer an attractive combination of accuracy, cost, and energy efficiency, and controlled experiments have demonstrated decimeter-to-centimeter positioning capability \cite{Wang2024SensorsSurvey,Zhu2024ComSTSurvey}. Moreover, optical positioning is naturally synergistic with optical wireless communications (OWC) and LiFi, enabling the co-existence of illumination, data, and localization within a unified infrastructure \cite{Haas2018LiFi,Kahn1997}.

Despite this promise, most practical VLP implementations still follow a multi-fixed-LED paradigm, where multiple ceiling luminaires act as anchors. This increases infrastructure density and deployment complexity, and it often requires calibration of beacon locations and/or channel parameters \cite{Zhuang2018Survey,Amsters2021RobotCalibration}. Experimental studies further highlight that real 3D VLP performance can be significantly impacted by receiver tilt and multipath reflections, emphasizing the need for robust and interpretable estimation methods \cite{Almadani2021OptCom}. At the same time, machine-learning-based VLP has gained significant attention as a means to compensate for nonlinearities and modeling errors \cite{TranHa2022MLReview}. While effective when abundant training data are available, such approaches may behave as black-box estimators and can require scenario- and hardware-specific retraining, offering limited physical transparency.

Motivated by the need to reduce infrastructure while preserving interpretability, this paper experimentally validates a model-based single beam-steered LED (SBL) architecture for 3D localization using a single photodiode receiver. Building on the ratio-based formulation in \cite{acuna2025model}, inter-orientation RSS ratios enable closed-form direction estimation without relying on data-driven fingerprinting. We implement a robotic 3D testbed using a commercial near-infrared LED on a two-axis gimbal and compare two closed-form estimators, generalized least squares (GLS) and weighted least squares (WLS), under a $K\!=\!9$ steering configuration. Beyond full 3D positioning, we evaluate direction-finding accuracy---a directly relevant metric for beam-steering-based OWC---and study the sensitivity of the estimates to the integration time, thereby characterizing the accuracy--latency trade-off of the proposed SBL method.

The remainder of this paper is organized as follows. Section~II presents the SBL system model and the GLS/WLS estimators. Section~III describes the 3D robotic testbed and experimental protocol. Section~IV reports the experimental results, including estimator comparison, direction-finding performance, and position-estimation stability. Section~V concludes the paper.











% =======================================================
\section{Introduction}
\label{sec:intro}
% =======================================================

\IEEEPARstart{I}{ndoor} positioning is a key enabler for emerging applications such as robotics, asset tracking, navigation, and context-aware wireless services. In this context, visible light positioning (VLP), and more broadly optical wireless positioning (OWP), have attracted significant interest because they can exploit LED-based infrastructure while offering high spatial reuse, low electromagnetic interference, and direct integration with optical wireless communication (OWC) systems. In particular, photodiode (PD)-based VLP remains attractive for practical deployment due to its low hardware complexity and high sampling capability.

Most RSS-based VLP systems, however, follow the conventional multi-anchor paradigm, where multiple fixed ceiling LEDs must be simultaneously visible to estimate the receiver position. Although effective, this approach increases infrastructure density, deployment cost, and calibration complexity, especially for full 3D localization. These limitations have motivated increasing interest in \emph{single-LED} positioning, where the goal is to recover the receiver state from only one optical anchor.

The literature shows that this goal is challenging because a single static optical measurement is not sufficient to uniquely determine a 3D position. As a result, existing single-LED approaches typically create additional information by increasing receiver-side complexity. Experimental 3D demonstrations based on a single LED have been reported using image sensors, gyroscopes, or hybrid off-the-shelf device combinations, while other single-LED systems employ multiple photodiodes with tilted geometries to recover angular diversity. More recently, single-LED single-PD schemes have also been explored, but often through tracking, receiver rotation, or machine-learning-based inference. Therefore, despite this progress, the experimental validation of a \emph{3D single-LED single-PD} positioning architecture with a fixed simple PD receiver and a fully model-based estimator remains very limited.

At the same time, machine-learning-based VLP has become an active research direction because it can compensate for nonlinearities and environmental distortions. Nevertheless, such methods are typically data-driven and scenario-dependent, and their learned mappings are less interpretable than estimators derived directly from radiometric geometry. By contrast, model-based approaches preserve a direct physical link between the optical channel, the measurement equations, and the recovered user state, which is particularly valuable when the goal is not only localization, but also system understanding and beam-control design. This point is especially relevant in the present context, since beam steering is increasingly regarded as a key enabler for high-capacity indoor OWC systems, making direction estimation and beam alignment natural companion functions of localization.

Motivated by this, our recent work in \cite{acuna2025model} introduced a ratio-based 3D positioning framework using a single beam-steered LED (SBL) and a single PD receiver. Under the ideal Lambertian model, inter-orientation RSS ratios cancel the receiver-orientation term and lead to closed-form direction estimators, including generalized least squares (GLS) and weighted least squares (WLS). Building on that formulation, this paper presents, to the best of our knowledge, the first experimental validation of 3D positioning using a single beam-steered LED and a single photodiode receiver. Using a robotic testbed and a $K=9$ steering configuration, we experimentally compare GLS and WLS for both 3D positioning and direction finding, and we further analyze the stability of the position estimate with respect to the integration time. In this way, the paper validates the practical feasibility of the SBL principle while retaining the main advantage of a model-based approach: an interpretable, closed-form, non-black-box localization framework suitable for beam-aware indoor optical wireless systems.